\documentclass{beamer}
\usepackage[utf8]{inputenc}

\usetheme{Madrid}
\usecolortheme{default}
\usepackage{amsmath,amssymb,amsthm}
\usepackage{txfonts}
\usepackage{tikz}
\usepackage{listings}
\usepackage{graphicx}
\usepackage{lmodern}

\setbeamertemplate{page number in head/foot}[totalframenumber]

\lstset{
    language=C,
    basicstyle=\ttfamily\small,
    keywordstyle=\color{blue},
    stringstyle=\color{orange},
    commentstyle=\color{green!60!black},
    numbers=left,
    numberstyle=\tiny\color{gray},
    breaklines=true,
    showstringspaces=false,
}


\title{5.3.2}
\author{Chiruvella Harshith Sharan - AI25BTECH11008}

\begin{document}

\frame{\titlepage}

%------------------------------------

\begin{frame}{Question}
The pair of linear equations
\[
\frac{3x}{2} + \frac{5y}{3} = 7 
\quad \text{and} \quad 
9x + 10y = 14
\]
is

\begin{enumerate}[label=(\alph*)]
\item consistent 
\item inconsistent 
\item consistent with one solution 
\item consistent with many solutions
\end{enumerate}
\end{frame}

%------------------------------------

\begin{frame}{Step 1: Clear Fractions}
Multiply the first equation by $6$:
\[
9x + 10y = 42 \tag{1}
\]
The second equation is
\[
9x + 10y = 14 \tag{2}
\]
\end{frame}

%------------------------------------

\begin{frame}{Step 2: Augmented Matrix}
The system can be written as:
\[
\begin{pmatrix}
9 & 10 & \vline & 42 \\
9 & 10 & \vline & 14
\end{pmatrix} \tag{3}
\]
Subtracting $R_1$ from $R_2$:
\[
\begin{pmatrix}
9 & 10 & \vline & 42 \\
0 & 0 & \vline & -28
\end{pmatrix} \tag{4}
\]
This implies $0x + 0y = -28$, which is a contradiction.
\end{frame}

%------------------------------------

\begin{frame}{Step 3: Conclusion}
Therefore, the two equations represent **parallel distinct lines**.

\[
\boxed{\text{The system is inconsistent.}}
\]
\end{frame}

%------------------------------------

\begin{frame}[fragile]{Python Code}
\begin{lstlisting}[language=Python]
import numpy as np
import matplotlib.pyplot as plt

# Define x-range
x = np.linspace(-5, 10, 400)

# Define the two lines: y = (constant - 9x)/10
y1 = (42 - 9*x) / 10   # from 9x + 10y = 42
y2 = (14 - 9*x) / 10   # from 9x + 10y = 14

# Plotting
plt.figure(figsize=(6,6))
plt.plot(x, y1, label=r'$9x+10y=42$', color='blue')
plt.plot(x, y2, label=r'$9x+10y=14$', color='red')
plt.xlabel('x')
plt.ylabel('y')
plt.axhline(0, color='black', linewidth=0.5)  # x-axis
plt.axvline(0, color='black', linewidth=0.5)  # y-axis
plt.title('Lines: $9x+10y=42$ and $9x+10y=14$ (Parallel, no intersection)')
plt.legend()
plt.grid(True, linestyle='--', linewidth=0.3)

plt.show()
\end{lstlisting}
\end{frame}

%------------------------------------

\begin{frame}[fragile]{C Code}
\begin{lstlisting}[language=C]
#include <stdio.h>

int main() {
    double a1=9, b1=10, c1=42;
    double a2=9, b2=10, c2=14;

    double r1 = a1/a2;
    double r2 = b1/b2;
    double r3 = c1/c2;

    if(r1 == r2 && r2 == r3)
        printf("Consistent: Infinite solutions\n");
    else if(r1 == r2 && r2 != r3)
        printf("Inconsistent system (No solution)\n");
    else
        printf("Consistent: Unique solution\n");

    return 0;
}
\end{lstlisting}
\end{frame}

%------------------------------------

\begin{frame}{Plot}
  \begin{figure}[h!]
    \centering
    \includegraphics[width=0.6\linewidth]{beamer/figs/fig1.png}
    \caption{Graph of the lines $9x+10y=42$ and $9x+10y=14$.}
    \label{fig:lines}
\end{figure}

\end{frame}

\end{document}

